% ============================================================
% Capitulo 3 -- Metodologia
% Tamanho-alvo: 8-12 paginas
% Ordem de escrita: PRIMEIRO (junto com Resultados).
% Mais denso em fatos verificaveis = menos espaco pra hand-waving.
% ============================================================

\chapter{Metodologia}
\label{cap:metodologia}

% ------------------------------------------------------------
\section{Vis\~ao geral do pipeline}
\label{sec:pipeline}

% ESCREVER (1 paragrafo): descricao breve dos blocos do pipeline.
% A figura \ref{fig:pipeline} (F0) abaixo e o diagrama em TikZ.

\begin{figure}[H]\centering
\resizebox{\linewidth}{!}{%
\begin{tikzpicture}[
    node distance=0.9cm and 0.7cm,
    box/.style    ={rectangle, rounded corners, draw, thick,
                    minimum width=2.0cm, minimum height=0.9cm,
                    align=center, font=\scriptsize},
    data/.style   ={box, fill=blue!10},
    proc/.style   ={box, fill=green!15},
    model/.style  ={box, fill=orange!18},
    eval/.style   ={box, fill=red!12},
    arr/.style    ={->, >=stealth, thick}
]
\node[data]  (raw)    {Dados brutos\\FNN, UPFD, Bluesky};
\node[proc, right=of raw]   (graph)  {Constru\c{c}\~ao\\de grafos};
\node[proc, right=of graph] (feat)   {\emph{Features}\\posicionais\\+\,BERT};
\node[model, right=of feat] (model)  {GNN\\GCN/GAT/SAGE};
\node[eval, right=of model] (val)    {Valida\c{c}\~ao\\k-fold pareado};
\node[eval, right=of val]   (out)    {M\'etricas\\F1, Cohen's $d$};

\draw[arr] (raw)   -- (graph);
\draw[arr] (graph) -- (feat);
\draw[arr] (feat)  -- (model);
\draw[arr] (model) -- (val);
\draw[arr] (val)   -- (out);
\end{tikzpicture}}
\caption{Vis\~ao geral do \emph{pipeline} de detec\c{c}\~ao. Dados brutos s\~ao
transformados em grafos de propaga\c{c}\~ao com \emph{features} estruturais
(e textuais via BERT, quando dispon\'iveis), passados por modelos GNN, e
avaliados com protocolo de k-fold pareado.}
\label{fig:pipeline}
\end{figure}

% ------------------------------------------------------------
\section{Datasets}
\label{sec:datasets}

% ESCREVER (1 paragrafo introdutorio + tabela T0 abaixo).
% Citacoes: \cite{shu2020fakenewsnet}, \cite{dou2021upfd},
% \cite{kleppmann2024atproto}, \cite{failla2024bluesky}.

\begin{table}[H]\centering
\caption{Datasets utilizados, com tamanho, fonte e papel no trabalho.}
\label{tab:datasets}
\small
\begin{tabular}{lrrll}
\toprule
\textbf{Dataset} & \textbf{N grafos} & \textbf{N classes} & \textbf{Fonte} & \textbf{Papel} \\
\midrule
FakeNewsNet PolitiFact   & $\sim$754   & 2 (50/50)   & \cite{shu2020fakenewsnet}      & Treino + valida\c{c}\~ao \\
UPFD-PolitiFact          & $\sim$314   & 2 (50/50)   & \cite{dou2021upfd}             & Cross-dataset \\
UPFD-GossipCop           & $\sim$5{.}464 & 2 (50/50) & \cite{dou2021upfd}             & Cross-dataset \\
Bluesky (acad\^emico)    & 168{.}463   & --          & \cite{failla2024bluesky}       & Aplica\c{c}\~ao (sem \emph{labels}) \\
\bottomrule
\end{tabular}
\end{table}

% Para cada dataset, 1 paragrafo: o que e, quem coletou, quantos grafos,
% feature dominante. Para Bluesky: declarar explicitamente que NAO tem
% ground-truth fake/real e e usado em modo demonstracao no Cap 5.

% ------------------------------------------------------------
\section{Constru\c{c}\~ao de grafos: FakeNewsNet}
\label{sec:construcao_fnn}

% ESCREVER (1.5 pagina): pipeline do 00_construir_grafos_fakenewsnet.py.
% Download CSVs -> BERT embedding dos titulos -> construcao grafo
% (raiz=noticia, filhos=tweets propagadores) -> features posicionais ->
% cache em data/fakenewsnet_posfull/.
% 3 variantes: posfull (768+3=771d), posmin (1d so is_root), posgrau (2d).

% ------------------------------------------------------------
\section{Folds compartilhados}
\label{sec:folds}

% ESCREVER (0.5 pagina): gerar_folds.py -- 10 folds estratificados,
% seed=42, salvos em folds_fnn.pt e reusados em todos os experimentos.
% Justificar: "Folds compartilhados sao pre-requisito para que
% scipy.stats.ttest_rel seja valido -- modelos comparados precisam ter
% sido avaliados nos *mesmos* exemplos."

% ------------------------------------------------------------
\section{Arquiteturas GNN e hiperpar\^ametros}
\label{sec:arquiteturas}

% ESCREVER (1 paragrafo introdutorio + tabela de hiperparametros abaixo).
% Defaults explicados na prosa.

\begin{table}[H]\centering
\caption{Hiperpar\^ametros usados nos modelos GNN. Defaults s\~ao iguais
entre arquiteturas; diferen\c{c}as anotadas na pr\'opria tabela.}
\label{tab:hyperparams}
\small
\begin{tabular}{lrrr}
\toprule
\textbf{Hiperpar\^ametro} & \textbf{GCN} & \textbf{GAT} & \textbf{SAGE} \\
\midrule
\texttt{hidden\_channels}     & 64     & 64     & 64 \\
\texttt{num\_layers}          & 3      & 3      & 3 \\
\texttt{dropout}              & 0{.}5  & 0{.}5  & 0{.}5 \\
\texttt{num\_heads} (atten\c{c}\~ao) & --     & 4      & -- \\
\texttt{pooling}              & \texttt{global\_mean\_pool} & idem & idem \\
\texttt{learning\_rate}       & 0{.}001 & 0{.}001 & 0{.}001 \\
\texttt{weight\_decay}        & $5\times10^{-4}$ & $5\times10^{-4}$ & $5\times10^{-4}$ \\
\texttt{batch\_size}          & 32     & 32     & 32 \\
\texttt{optimizer}            & Adam   & Adam   & Adam \\
\texttt{scheduler}            & \multicolumn{3}{c}{ReduceLROnPlateau (factor=0{.}5, patience=5)} \\
\texttt{early\_stopping}      & \multicolumn{3}{c}{patience=7 sobre F1-macro de valida\c{c}\~ao} \\
\texttt{max\_epochs}          & 30     & 30     & 30 \\
\texttt{seed}                 & 42 (com varredura para 10 \emph{seeds} em \S\ref{sec:topo_sem_texto}) \\
\bottomrule
\end{tabular}
\end{table}

% ------------------------------------------------------------
\section{\emph{Baselines} n\~ao-GNN}
\label{sec:baselines}

% ESCREVER (1 pagina):
% LogReg-BERT-root: feature = x[0] do grafo (BERT do titulo da raiz).
% RF-tabular: features = [num_nodes, grau_root].
% Justificar: tetos inferiores + diagnostico de confound. Se RF-tabular
% sozinho atinge F1 alto, GNN nao esta capturando nada alem de tamanho.

% ------------------------------------------------------------
\section{Protocolo de avalia\c{c}\~ao}
\label{sec:protocolo}

% ESCREVER (1 pagina):
% F1-macro como metrica primaria (justificar -- classes balanceadas no
% FNN/UPFD; evita inflacao por classe majoritaria).
% k-fold pareado + scipy.stats.ttest_rel para comparacao de modelos.
% Cohen's d para analise estrutural (\S 4.5).
% Reprodutibilidade caveat: "Resultados foram obtidos em CPU (Python 3.12.10,
% torch 2.10.0+cpu); seeds sao fixadas mas determinismo nao e bit-exact
% entre maquinas -- variacao tipica na 3a casa decimal."

% ------------------------------------------------------------
\section{An\'alise estratificada do erro}
\label{sec:metod_estratificacao}

% ARGUMENTO: F1 agregado mede o "quanto" -- nao mede o "onde". Esta secao
% descreve tres tecnicas complementares (LOWESS, percentil empirico, bootstrap
% pareado) que cruzam metrica estrutural com acerto do modelo, identificando
% onde dentro da distribuicao o modelo brilha ou falha. Cada tecnica e
% defendida pela limitacao especifica do metodo que ela resolveria se usada
% em isolamento.

% RESUMO DO QUE VAI AQUI (~0.7 pagina, sem figuras):
%   1. LOWESS (smoother por janela movel sobre dados ordenados): visualiza
%      continuamente P(acerto | metrica), sem cherry-picking de bins.
%      Implementacao via rolling-mean de janela proporcional (10% da amostra,
%      minimo 30) para evitar dependencia do statsmodels.
%   2. Percentil empirico (vs z-score): cascatas em redes sociais sao
%      heavy-tail; usar limiar de outlier baseado em sigma e mal-definido.
%      Reportamos os 6 subgrupos da Tabela \ref{tab:outliers} via p5/p95
%      empiricos.
%   3. Bootstrap pareado para CI 95% da diferenca F1(SAGE) - F1(RF) em
%      subsets (n_boot=2000). Justifica afirmacoes do tipo "ganho do SAGE
%      e significativo neste subset" mesmo onde nao temos folds (subsets
%      definidos post-hoc no test).
% Fonte: scripts 28, 29 e 30 do mega research (Fase 9 do PIPELINE.md).

% ESCREVER (Cesar): justificar por que cada tecnica e a "minima necessaria"
% pra responder a pergunta dela. Citacao opcional: Cleveland 1979 (LOWESS),
% Cohen 1988 (effect size), Efron & Tibshirani 1993 (bootstrap).

% ------------------------------------------------------------
\section{An\'alise de explicabilidade}
\label{sec:analise_explicabilidade}

% ESCREVER (0.5 pagina): GNNExplainer aplicado em SAGE com features
% [is_root, grau_norm], 20 amostras estratificadas por (classe x tamanho)
% no UPFD-GossipCop test, analise de hop-distance (1-hop, 2-hop, 3+).

% ------------------------------------------------------------
\section{Classificador dual: regra de combina\c{c}\~ao \emph{post-hoc}}
\label{sec:dual}

% ESCREVER (0.5 pagina) -- CRITICO: introduzir como heuristica post-hoc,
% NAO componente metodologico a priori.
%
% Texto sugerido:
% "Para a aplicacao pratica (Capitulo 5), combinamos o classificador textual
% (LogReg-BERT) e o topologico (RF estrutural) por uma regra de combinacao.
% Essa regra nao e uma decisao metodologica a priori; e uma heuristica
% derivada APOS observar os resultados de \S\ref{sec:concordancia}.
% Sua justificativa empirica e a discussao das alternativas estao naquela
% secao."
%
% NAO detalhar a regra aqui. So anunciar.

% ------------------------------------------------------------
\section{Reprodutibilidade}
\label{sec:reprodutibilidade}

% ESCREVER (0.5 pagina): requirements.txt pinado, seed=42 universal,
% README.md como fonte autoritativa. Listar caveat de bit-exact
% (CUDA/cuDNN nao-determinismo). Ja esta no README -- so citar.
